\subsection{Energy conditions}\label{sec:energy}

Let $\GH_{\mathrm{DEC}}$ be the full subcategory of spacetimes whose Einstein
tensor satisfies the dominant energy condition. It is closed under restriction
to causally convex open subsets and contains $\Mink$.

\begin{proposition}[The auxiliary spacetime is unavailable]\label{prop:dec}
If $M\in\GH_{\mathrm{DEC}}$ is flat on a neighbourhood of some totally geodesic
Cauchy surface, then $M$ is flat. In particular the spacetime $M'$ of Appendix~\ref{sec:proofs} lies in $\GH_{\mathrm{DEC}}$ only if
$g$ is flat on $T_{2r}$.
\end{proposition}

\begin{proof}
On the Cauchy surface $S$ the metric is flat and the Einstein tensor vanishes. By the conservation theorem
for matter obeying the dominant energy condition \citep[\S4.3]{HE1973}, the
Einstein tensor vanishes on its domain of dependence, i.e.\ everywhere. In the
form we use: let $t$ be a Cauchy temporal function with $S=t^{-1}(0)$
\citep[Thm.~1.2]{BS2006} and $p\in J^+(S)$. The set $J^-(p)\cap J^+(S)$ is compact, and
it is the union of the segments of integral curves of $\nabla t$ that run from
$S$ to $\partial J^-(p)$ (a future-directed curve that leaves the past set
$J^-(p)$ does not return). Minguzzi's form of the theorem \citep[Thm.~1]{Minguzzi2015},
which needs no time function, then gives $G=0$ at $p$ from $G=0$ on $S$;
$J^-(S)$ is the time reverse. For
$M'$, any slice $\{t=c\}$ with $c<-3r$ is totally geodesic and has a flat
neighbourhood.
A vacuum globally hyperbolic spacetime with flat data $(\Sigma,h,0)$ on a
Cauchy surface is a development of those data and so embeds isometrically in
their maximal development \citep[Thm.~3]{CBG1969}. That maximal development is
the flat static spacetime $\mathbb R\times\Sigma$, $-dt^2+h$ (not Minkowski in
general; e.g.\ $\Sigma=T^3$): it is a development, and if a development
properly extended it, a point $p$ outside its image would lie on a timelike
geodesic from $\Sigma$ whose image under the extension coincides with the
corresponding complete geodesic of $\mathbb R\times\Sigma$, so $p$ would lie in
the image. Hence the spacetime is flat.
\end{proof}

The same reasoning constrains any repair that keeps the original spacetime and
inserts small flat pockets near the points $w$, $v$ and the detour of the proof of Theorem~\ref{thm:main}: in a vacuum region
whose causal past on some Cauchy surface is vacuum and unmodified, the
conservation theorem forces the modified spacetime to be vacuum there with the
same data, so no such pocket can be inserted there. Pockets must be fed by matter brought
in from the past and routed around $D(\Sigma)$ afterwards. The proof of
Theorem~\ref{thm:main} goes through verbatim in $\GH_{\mathrm{DEC}}$ given
the following sufficient condition, which we leave open.

\begin{quote}
\textbf{Pocket problem.} Given $M\in\GH_{\mathrm{DEC}}$ and $\Sigma$ as in Lemma~\ref{lem:D},
find $M''\in\GH_{\mathrm{DEC}}$ containing an isometric copy of a neighbourhood
of $D_M(\Sigma)$ whose image is a domain of dependence, and exactly flat
causally convex regions containing points $w\ll x$ and $v\gg b$ joined by the
flat domain of dependence of a slab.
\end{quote}

The constraint is an inequality, not a field equation, so constructions can be
kinematic, and explicit models with flat regions chronologically adjacent to
curved dominant-energy regions exist (Vaidya spacetimes before an incoming, or
after an outgoing, null-dust shell). Thin shells will not suffice in general,
since null junctions require matching intrinsic geometry. A solution would
give Theorem~\ref{thm:main} for spacetimes with physically reasonable matter;
a vacuum version would require bounded interior flat regions adjacent to
arbitrary curved vacuum regions, which we have not found in the gluing
literature (an abstract-level survey: Carlotto--Schoen, Corvino--Schoen,
Aretakis--Czimek--Rodnianski, Czimek--Rodnianski).
